%%%%%%%%%%%%%%%%%%%%%%%%%%%%%%%%%%%%%%%%%%%%%%%%%%%%%%%%%%%%%
%  CHAPTER 1: FOUNDATIONS
%  Template for first chapter
%%%%%%%%%%%%%%%%%%%%%%%%%%%%%%%%%%%%%%%%%%%%%%%%%%%%%%%%%%%%%

\chapter{Some Mathematical Preliminaries}
\label{ch:foundations}

\section{Introduction}

This chapter introduces some fundamental mathematics that will be needed in the later chapters. The terms appearing in this chapter should be familiar to any student of physics. However, due to the intuitive way these terms are usually presented in any first course of calculus at the university freshman level, students are not always familiar with the more "mathematical" way of presenting these concepts. Since in the later study of basic differential geometry these will work as a necessary prerequisite vocabulary, it is better to be familiar with the language a bit. 









\section{Revisiting Calculus Notions}

Let $ U \subseteq \mathbb{R}^n$ and $ f: U \longrightarrow \mathbb{R}$. We define the partial derivatives as follows:—

\begin{align*}
    & \pp{f}{x^i} \bigg|_a := \displaystyle\lim_{h \to 0} \frac{ f \lera{ a^1, \ldots, a^i + h, \ldots, a^n} - f \lera{ a^1, \ldots, a^i, \ldots, a^n}}{h} \\
\end{align*}

If $ \pp{f}{x^i} \big|_a$ exists for all $ i$ $ \lera{ 1 \leq i \leq n}$ for each point $ a \in U$, then we have n new functions in U. If these functions $ \pp{f}{x^i}$ are also continuous on U, then we say that f is continuosly differentiable on U, i.e, $ f \in C^1(U)$.

Existence of partial derivatives at point a is not enough to guarantee the differentiability of f at a. To illustrate this point, we use an example. Consider the function, defined on $ \mathbb{R}^2$, by

\begin{align*}
    \begin{cases}
        f(x,y) = \frac{xy}{x^2 + y^2} \\
        f(0,0) = 0
    \end{cases} \\
\end{align*}

We evalute the following,

\begin{align*}
    \pp{f}{x} \bigg|_{(0,0)} &= \displaystyle\lim_{h \to 0} \frac{f(h,0) - f(0,0)}{h} \\
    &= \displaystyle\lim_{h \to 0} \frac{0-0}{h} = 0 \\
\end{align*}

Similarly, $ \pp{f}{y} \big|_{(0,0)} = 0$

So, both of the partial derivaties exist at the point $ (0,0)$. But the function is still not differentiable. To prove that, it suffices to show that the function is not continuous at $ (0,0)$.

To decide for continuity, we first try to find out the limit of $ f(x,y)$ at $ (0,0)$. Considering the two different paths as in the picture

\begin{center}
\begin{tikzpicture}[>=stealth, scale=1.1]
    \draw[->] (-2,0) -- (2.2,0) node[right] {$ x$};
    \draw[->] (0,-1) -- (0,2) node[above] {$ y$};
    \draw (-1.4,-1.4) -- (1.7,1.7) node[right] {$ y=x$};
    \draw[->] (1.0,0) -- (0.7,0);
    \node[below] at (0.85,-0.05) {I};
    \draw[->] (0.95,0.95) -- (0.7,0.7);
    \node[above left] at (1.15,0.85) {II};
\end{tikzpicture}
\end{center}

\begin{align*}
    \displaystyle\lim_{\substack{(x,y) \to (0,0) \\ \text{along path II}}} f(x,y) &= \displaystyle\lim_{x \to 0} \frac{x \cdot x}{x^2 + x^2} \\
    &= \displaystyle\lim_{x \to 0} \frac{x^2}{2x^2} \\
    &= \frac{1}{2} \\
\end{align*}

\begin{align*}
    & \displaystyle\lim_{\substack{(x,y) \to (0,0) \\ \text{along I}}} f(x,y) = 0 \\
\end{align*}

$ \therefore$ limit along I $ \neq$ limit along II

$ \therefore$ The limit DNE at the point $ (0,0)$ and so $ f(x,y)$ is not continuos at $ (0,0)$. Hence it is not differentiable at the origin even though both of the partial derivatives exist.

Let, $ f: U \subseteq \mathbb{R}^n \longrightarrow \mathbb{R}$

If f has continuous first partial derivatives throughout U, then $ f \in C^1(U)$.

\begin{definition}
    We consider $ f \in C^r(U)$ iff for all $ j \leq r$,

    \begin{align*}
        & \frac{\partial^j f}{\partial x^{i_1} \partial x^{i_2} \ldots \partial x^{i_j}} \bigg|_a \qquad \lera{ 1 \leq i_1, \ldots, i_j \leq n} \text{ exists} \\
    \end{align*}

    for all $ a \in U$ and they are continuos on U.
\end{definition}

Here the order of the derivatives doesnt matter.

A function $ f: U \subseteq \mathbb{R}^n \longrightarrow \mathbb{R}$ is said to be differentiable at a point a, iff there exists a linear expression $ \sum_{i=1}^{n} b_i \lera{ x^i - a^i}$ which satisfies the following

\begin{align*}
    & \displaystyle\lim_{x \to a} \frac{ f(x) - f(a) - \sum_{i=1}^{n} b_i \lera{ x^i - a^i}}{\| x-a \|} = 0 \\
\end{align*}

So, we can say that $ f(a) + \sum_{i=1}^{n} b_i \lera{ x^i - a^i}$ approximates $ f(x)$ near the point a.

Using the above expression, or, you can say keeping it at the back of our mind, we can come to another def\textsuperscript{n} of differentiability. Here, a function is said to be differentiable at a, if there exists constants $ b_1, \ldots, b_n$ and a function $ r(x,a)$ such that

\begin{align*}
    & f(x) = f(a) + \sum_{i=1}^{n} b_i \lera{ x^i - a^i} + \| x-a \| r(x,a) \\
    & \text{and } \displaystyle\lim_{x \to a} r(x,a) = 0 \\
\end{align*}

When a function $ f: U \subseteq \mathbb{R}^n \longrightarrow \mathbb{R}$ is differentiable at a, then f is continuous at a and $ b_i = \pp{f}{x^i} \big|_a$

\begin{align*}
    & \text{So, } f(x) = f(a) + \sum_{i=1}^{n} \pp{f}{x^i} \bigg|_a \lera{ x^i - a^i} + r(x,a) \| x-a \| \\
\end{align*}

Keep in mind that, f is differentiable at a, merely guarantees the existence of $ \pp{f}{x^i} \big|_a$. It doesnt guarantee that f is $ C^1$ at the point a.

If f is $ C^1$ in a neighborhood of a, then f is guaranteed to be differentiable at that neighbourhood.

Thus if $ f \in C^1(U)$ where U is an open set, then f is differentiable on U.

Let's generalize the notions learned so far.

Assume that, $ f: U \subseteq \mathbb{R}^n \longrightarrow V \subseteq \mathbb{R}^m$

The first thing that we must make clear is as follows:—

\begin{enumerate}
    \item $ f(x) = \lera{ f^1(x), f^2(x), \ldots, f^m(x)}$

    we call the individual functions $ f^i$ to be coordinate functions. We define them as follows. Suppose $ \pi^i : \mathbb{R}^m \longrightarrow \mathbb{R}$ is a function given by $ \pi^i \lera{ x^1, x^2, \ldots, x^m} = x^i$.

    Now, $ f^i := \pi^i \circ f : U \subseteq \mathbb{R}^n \longrightarrow \mathbb{R}$

    So it's clear from the above that,

    \begin{align*}
        & f(x) = \lera{ f^1(x), \ldots, f^m(x)} \\
    \end{align*}

    \item f is continuous on U if $ f^1, \ldots, f^m$ are continuous on U.

    f is differentiable on U if $ f^1, \ldots, f^m$ are

    f is $ C^r(U)$ if $ f^1, \ldots, f^m$ are $ C^r(U)$
\end{enumerate}

For ease of computation, we assume that the points of $ \mathbb{R}^n$ or $ \mathbb{R}^m$ are given by column matrices. So, for instance,

\begin{align*}
    & a \in \mathbb{R}^n \text{ and } a = \begin{bmatrix} a^1 \\ a^2 \\ \vdots \\ a^n \end{bmatrix} \\
\end{align*}

In this case

\begin{align*}
    & f(x) = \begin{bmatrix} f^1(x) \\ f^2(x) \\ f^3(x) \\ \vdots \\ f^m(x) \end{bmatrix} \\
\end{align*}

We give another def\textsuperscript{n} for differentiability.

$ f: U \subseteq \mathbb{R}^n \longrightarrow \mathbb{R}^n$ is differentiable at a point $ a \in U$ if there exists a matrix A and a function $ r(x,a)$ that satisfies the following two conditions:—

\begin{align*}
    & f(x) = f(a) + A \lera{ x-a} + r(x,a) \| x-a \| \\
    & \text{and } \displaystyle\lim_{x \to a} r(x,a) = 0 \\
\end{align*}

If f is differentiable at the point a, then A is the unique Jacobian matrix at a.

Now comes the def\textsuperscript{n} of the Jacobian matrix.

For any point $ p \in U$, we define the following:—

\begin{align*}
    & \text{Jac} f(p) = \frac{\partial \lera{ f^1, \ldots, f^m}}{\partial \lera{ x^1, \ldots, x^n}} \bigg|_p :=
    \begin{bmatrix}
        \pp{f^1}{x^1} \big|_p & \pp{f^1}{x^2} \big|_p & \cdots & \pp{f^1}{x^n} \big|_p \\[2ex]
        \pp{f^2}{x^1} \big|_p & \pp{f^2}{x^2} \big|_p & \cdots & \pp{f^2}{x^n} \big|_p \\[2ex]
        \vdots & & \ddots & \vdots \\[2ex]
        \pp{f^m}{x^1} \big|_p & \cdots & \cdots & \pp{f^m}{x^n} \big|_p
    \end{bmatrix} \\
\end{align*}

This is the Jacobian matrix of the function f at the point P. This matrix defines a map $ Df \big|_p : U \longrightarrow \mathbb{R}^m$. At each point of U, we get different different maps.

So, for the function f (when it is differentiable at the point a),

\begin{align*}
    & f(x) = f(a) + \lera{ Df \big|_a} \lera{ x-a} + r(x,a) \| x-a \| \\
\end{align*}

$ Df \big|_a$ acting on $ (x-a)$ which basially is

\begin{align*}
    & \text{Jac} f(a) \begin{bmatrix} x^1 - a^1 \\ x^2 - a^2 \\ \vdots \\ x^n - a^n \end{bmatrix} \\
\end{align*}

Now, we come to prove chain rule in its most general version.

\begin{align*}
    & f: U \subseteq \mathbb{R}^n \longrightarrow V \subseteq \mathbb{R}^m \\
    & g: V \subseteq \mathbb{R}^m \longrightarrow \mathbb{R}^k \\
    & \therefore g \circ f : U \subseteq \mathbb{R}^n \longrightarrow \mathbb{R}^k \\
\end{align*}

The chain rule says,

\begin{align*}
    & D g \circ f \big|_p = Dg \big|_{f(p)} \circ Df \big|_p \\
\end{align*}

Using the general case of the chain rule, we can derive important special cases:—

\begin{enumerate}
    \item
    \begin{align*}
        & f: U \subseteq \mathbb{R}^n \longrightarrow V \subseteq \mathbb{R}^m \\
        & g: V \subseteq \mathbb{R}^m \longrightarrow \mathbb{R}^k \\
        & g \circ f : U \subseteq \mathbb{R}^n \longrightarrow \mathbb{R}^k \\
    \end{align*}

    \begin{align*}
        & \frac{\partial \lera{ g \circ f}^i}{\partial x^j} \bigg|_p = \pp{g^i}{u^k} \bigg|_{f(p)} \pp{f^k}{x^j} \bigg|_p \\
    \end{align*}

    \item
    \begin{align*}
        & f: U \subseteq \mathbb{R}^n \longrightarrow V \subseteq \mathbb{R}^m \\
        & g: V \subseteq \mathbb{R}^m \longrightarrow \mathbb{R} \\
        & \therefore g \circ f : U \subseteq \mathbb{R}^n \longrightarrow \mathbb{R} \\
    \end{align*}

    \begin{align*}
        & \frac{\partial \lera{ g \circ f}}{\partial x^i} \bigg|_p = \pp{g}{u^j} \bigg|_{f(p)} \pp{f^j}{x^i} \bigg|_p \\
    \end{align*}

    \item
    \begin{align*}
        & f: U \subseteq \mathbb{R}^n \longrightarrow V \subseteq \mathbb{R} \\
        & g: V \subseteq \mathbb{R} \longrightarrow \mathbb{R} \\
        & \therefore g \circ f : U \subseteq \mathbb{R}^n \longrightarrow \mathbb{R} \\
    \end{align*}

    \begin{align*}
        & \frac{\partial \lera{ g \circ f}}{\partial x^i} \bigg|_p = \pp{g}{u} \bigg|_{f(p)} \pp{f}{x^i} \bigg|_p \\
    \end{align*}

    \item
    \begin{align*}
        & f: U \subseteq \mathbb{R} \longrightarrow V \subseteq \mathbb{R} \\
        & g: V \subseteq \mathbb{R} \longrightarrow \mathbb{R} \\
        & g \circ f : U \subseteq \mathbb{R} \longrightarrow \mathbb{R} \\
    \end{align*}

    \begin{align*}
        & \frac{d \lera{ g \circ f}}{d x} \bigg|_p = \dd{g}{u} \bigg|_{f(p)} \dd{f}{x} \bigg|_p \\
    \end{align*}

    \item
    \begin{align*}
        & f: U \subseteq \mathbb{R} \longrightarrow V \subseteq \mathbb{R}^n \\
        & g: V \subseteq \mathbb{R}^n \longrightarrow \mathbb{R}^m \\
        & g \circ f : U \subseteq \mathbb{R} \longrightarrow \mathbb{R}^m \\
    \end{align*}

    \begin{align*}
        & \frac{\partial \lera{ g \circ f}^i}{\partial x} \bigg|_p = \pp{g^i}{u^k} \bigg|_{f(p)} \pp{f^k}{x} \bigg|_p \\
    \end{align*}

    \item
    \begin{align*}
        & f: U \subseteq \mathbb{R} \longrightarrow V \subseteq \mathbb{R}^m \\
        & g: V \subseteq \mathbb{R}^n \longrightarrow \mathbb{R} \\
        & g \circ f : U \subseteq \mathbb{R} \longrightarrow \mathbb{R} \\
    \end{align*}

    \begin{align*}
        & \therefore \frac{\partial \lera{ g \circ f}}{\partial x} \bigg|_p = \pp{g}{u^i} \bigg|_{f(p)} \pp{f^i}{x} \bigg|_p \\
    \end{align*}
\end{enumerate}

If f and g are of class $ C^r$ (or smooth) on U and V, then $ h = g \circ f$ is of class $ C^r$ (smooth) on U.